\section{Counterarguments and the conditional conclusion}

The first counterargument is that frontier providers will do this themselves. They already publish agent toolkits, browser interfaces and engineering reports about harness design \citep{anthropiceffective,openaiagents}. Their distribution and model access are formidable. The aftermarket case rests on a different comparative advantage: narrow domain knowledge, customer-specific context, permitted integrations, acceptance rules and exception operations that a general consumer product may not prioritize. That advantage must be measured after native product upgrades. If it vanishes, the specialist should pivot or stop, not rename ordinary prompting as a moat.

The second counterargument is simpler: deterministic software is better. For a stable API and a known procedure, this may be true. A conventional integration is often easier to audit, cheaper per call, and less variable. The model earns its place where instructions, documents, interfaces or exceptions vary enough that rigid rules fail. An effective harness can use deterministic software for the stable parts and a model for the ambiguous parts. The evaluation must be willing to conclude that the model contributed no incremental value to a given workflow.

The third counterargument is that the nudge tax is just the price of necessary human judgment. In some tasks it is. A person may need to decide, approve, identify themselves or review a high-consequence outcome. Our target is the tactical work of repeatedly reconstructing context and steering routine steps. A harness that shifts this work to an opaque review queue, then calls it automation, has not solved the problem. The user-active minutes, waiting time and quality of the final result will reveal the difference.

The fourth counterargument is empirical. Benchmarks of agent action and benign over-refusal do not establish current consumer-product demand, refusal prevalence, or a profitable mass market. We agree. WebArena, TheAgentCompany and $\tau$-bench demonstrate useful ways to specify and verify multi-step outcomes \citep{webarena,agentcompany,taubench}; XSTest and OR-Bench demonstrate calibrated refusal testing \citep{xstest,orbench}. None ranks the commercial opportunities open to non-frontier laboratories. A market claim needs paired customer tasks, a credible alternative, and task-level economics. Even a productivity study of assisted customer support does not automatically transfer to autonomous cross-service completion \citep{genaiwork}.

Our conclusion is therefore forceful but conditional. The commercially important frontier is often between a model that can reason about a task and a product that can deliver an accepted result. For a team unable to improve the frontier model itself, an aftermarket harness offers a concrete research lever: change the surrounding task contract, context, control and verifier while holding intelligence fixed. The value proposition is to remove repeated tactical nudging from legitimate work, not to turn every refusal into compliance. The saleable unit is a verified completion at a known total cost.

Whether this is \emph{the greatest} non-frontier opportunity remains open. The thesis survives only if a specialist can repeatedly recover valuable permitted tasks, beat skilled nudging and conventional software after full costs, and retain an advantage as provider products improve. A research program that can fail those tests is more useful than a list of impressive demos. It can tell founders where to build, researchers what to measure, and customers what finished work they are actually buying.
